% ECONOMICS_PROBLEM: {"id": "G14-279", "title": "AI trading and market efficiency", "jel_code": "G14", "source_id": "OP-0319"}
% !TeX program = xelatex
\documentclass[11pt,a4paper]{article}
\usepackage[margin=23mm]{geometry}
\usepackage{fontspec}
\setmainfont{Latin Modern Roman}
\usepackage{amsmath,amssymb,mathtools}
\usepackage{xcolor,enumitem,booktabs,longtable,array,xurl,hyperref}
\definecolor{muted}{HTML}{53636C}
\hypersetup{colorlinks=true,urlcolor=blue,linkcolor=blue}
\setlength{\parindent}{0pt}
\setlength{\parskip}{5pt}
\newcommand{\R}{\mathbb R}
\newcommand{\E}{\mathbb E}
\newcommand{\N}{\mathbb N}
\newcommand{\ind}{\mathbf 1}
\newcommand{\Var}{\operatorname{Var}}
\newcommand{\Cov}{\operatorname{Cov}}
\newcommand{\supp}{\operatorname{supp}}
\DeclareMathOperator*{\argmin}{arg\,min}
\DeclareMathOperator*{\argmax}{arg\,max}
\newcommand{\entrymeta}[3]{{\sffamily\small\color{muted}\textbf{\texttt{#1}}\quad|\quad #2\par #3\par}}
\newcommand{\Problemref}[1]{\texttt{#1}}
\newcommand{\JELref}[2]{\texttt{#2} (JEL \texttt{#1})}
\begin{document}
\section*{AI trading and market efficiency}
\textbf{Problem ID:} \texttt{G14-279}\quad \textbf{JEL:} \texttt{G14}\par
% BEGIN ECONOMICS STATEMENT
\label{problem:OP-0319}
{\sffamily\small\color{muted}Original subject group: Asset pricing and financial markets\par}
\label{definitions:problem:30007601}
\textbf{Audit classification:} Published research agenda. The agenda supports studying algorithmic trading and stability. A specific learning game and collusion statistic are editorial; strategic algorithmic collusion is outside a strictly nongame catalogue.
\subsection*{Definitions and precise research target}
Fix a repeated asset market with a fundamental payoff \(V\), price \(P\), private signals, an order-processing rule and strategic traders. A learning system is specified by its training observations, objective, action space, update rule and computational restrictions. Let \(s\) denote the market share using a specified family of artificial-intelligence trading systems. Define an intervention that varies \(s\), with competitive and informational responses of other traders allowed. At share \(s\), measure price inefficiency \(E[(P-V^{\mathrm{fund}})^2]\), where \(V^{\mathrm{fund}}\) is the declared date-specific fundamental present value, effective spreads, trading volume and the expected resource loss of market disruption. Define a collusion statistic as the persistent excess of joint trading profits or spreads over a separately computed competitive benchmark under identical information and technology. The task is to identify how these outcomes change with \(s\) and characterize conditions under which learning supports tacit collusion or correlated destabilizing behavior. Vary signal precision, feedback, algorithm diversity and entry while keeping the comparison explicit. Historical backtests cannot reveal all equilibrium responses to widespread adoption; live or experimentally controlled validation must address adaptation and data leakage. Particular learning environments can exhibit inefficient outcomes without proving that all AI trading has that effect.

\textbf{Additional formal definitions and scope (editorial).} Choose a finite evaluation horizon \(T\), discount weights and a terminal payoff \(V\in L^2\), and fix the market's auction/clearing rule and admissible orders. Fundamental value at date \(t\) is a prespecified conditional present value, rather than the undiscounted terminal payoff at all dates. The algorithm specification includes initialization, exploration randomness, training/evaluation split and the optimization target. The share \(s\) is explicitly measured by capital, trader count or order flow. Compare price error and profits with an equilibrium benchmark using identical signals, risk tolerance, action space and transaction costs. Excess profits alone do not identify collusion: distinguish market power, risk-bearing rents and information rents, and define the tested conduct restriction. Strategic collusion variants use game-theoretic concepts; a catalogue restricted to nongame economic theory should keep an adoption/identification formulation and label the strategic extension as outside that strict boundary.

For the identification part, declare an admissible class \(\Theta\) of structural models, including preferences, technologies, shock laws, measurement/sampling rules and any equilibrium selector. If \(O\) denotes the complete observable history and \(T(\theta)\) the target defined above, the identified set is \(\mathcal I_T=\{T(\theta):\theta\in\Theta,\ \mathcal L_\theta(O)=\mathcal L_{\rm obs}(O)\}\); point identification means this set is a singleton. A \(do(a)\) comparison replaces stated policy/technology equations while fixing the declared exogenous law and re-solving the remaining equations. These definitions do not assert that an identification theorem holds without further restrictions.
\subsection*{Variants and assumption changes}
\begin{enumerate}
\item \textbf{Competitive learning adoption (editorial).} Treat traders as price takers and specify information acquisition; estimate adoption effects on price information and spreads.
\item \textbf{Strategic algorithms (editorial; game-theoretic).} Fix a repeated trading game and learning rule; test explicitly defined departures from a competitive conduct benchmark, permitting correlated algorithm behavior.
\end{enumerate}
\subsection*{References and source audit}
\begin{enumerate}
\item Official January 9, 2026 web version; locator: The Financial System / Technological innovation; Prudential Architecture / Resilience. Broad agenda; exact model editorial. URL: \url{https://www.bankofengland.co.uk/research/bank-of-england-agenda-for-research}.
\end{enumerate}
\begin{enumerate}\raggedright
\item\hypertarget{definitions:source:30007601:1}{} Bank of England. 2026. \emph{Bank of England Agenda for Research, 2025–2028}. The Financial System / Technological innovation; Prudential Architecture / Resilience.
\url{https://www.bankofengland.co.uk/research/bank-of-england-agenda-for-research}.
\end{enumerate}

\subsection*{Common conventions: Source mathematical conventions}

These conventions apply to each entry unless explicitly replaced.

\textbf{Models and identification.} A primitive model \(m\) specifies all probability laws, feasible choices, information, equilibrium and measurement rules used by its target. The admissible class \(\mathcal M\) is fixed independently of outcomes. If \(P_O(m)\) is its observable probability law and \(\tau(m)\) the target, its identified set at observable law \(P\) is
\[
 \mathcal I_\tau(P)=\{\tau(m):m\in\mathcal M,\ P_O(m)=P\}.
\]
Point identification means that this set is a singleton. An empty set signals incompatibility of data and assumptions. Coordinate bounds need not describe the joint set. Bounds are sharp only if every retained target is attainable or arbitrarily approachable by compatible models. Finite bounds require explicit restrictions on unobserved quantities; unrestricted confounding, hidden wealth, extrapolation or arbitrary equilibrium selection can yield uninformative sets.

\textbf{Causal interventions.} A potential outcome \(Y_i(p)\) is the outcome under a complete policy regime \(p\), including timing, financing, information and market spillovers. The target population and units are fixed before treatment. With interference, use the complete assignment vector or a justified exposure mapping; individual no-interference notation alone cannot identify an economy-wide intervention. A difference of mean potential outcomes does not require the cross-world joint distribution. Conditioning on post-treatment migration, survival, enrollment or employment changes the estimand and needs separate justification.

\textbf{Statistical statements.} An estimator, confidence set, forecast or test specifies a sampling law, model class and loss/coverage criterion. Uniform \((1-\alpha)\)-coverage means \(\inf_{m\in\mathcal M}\Pr_m\{\tau(m)\in C_n\}\geq1-\alpha\), or an explicitly asymptotic version. The whole parameter space trivially has coverage, so informativeness requires a stated width, risk or power objective. A forecasting improvement is relative to a named benchmark, information set and evaluation population.

\textbf{Equilibrium and optimization.} An equilibrium comprises feasible plans, optimality, beliefs where needed, budget constraints and all market/resource clearing. The primitives or a stated benchmark define each correspondence \(\mathcal E(p;m)\). Nonemptiness is assumed or proved, never inferred just from compact policy sets. A selected-equilibrium target uses a declared selection rule; a robust target quantifies over all permitted equilibria. Use suprema or infima unless attainment is established. An \(\varepsilon\)-optimal policy is feasible and within \(\varepsilon\) of the finite optimum value. Report an empty feasible set separately.

\textbf{Welfare and accounting.} Normative utility, interpersonal normalization, social weights, numeraire, discounting and the use of public receipts are primitives. Behavioral utility need not coincide with the welfare criterion. Household welfare includes ownership income and rents once; adding profits again to compensating variation generally double-counts them. A monetary equivalent is defined by an explicit transfer and price convention, with monotonicity and range conditions. Average effects, mechanism contributions, transfers and welfare losses are different targets. Infinite sums require convergence and dynamic feasible plans require terminal or no-Ponzi conditions.

\textbf{Source versions.} A working-paper date, revision date and journal publication year are distinguished. A DOI for a journal version is not automatically the DOI of an earlier manuscript. Locators refer to the stated version and printed pages where specified. A metadata-only check does not verify a claim about a full-text conclusion. Every entry's source subsection narrows the provenance claim accordingly.
% END ECONOMICS STATEMENT
\end{document}
